%
% nichtungerade.tex -- Graph der Funktion f(x) = (x+1)/(1+x^2)
%
% (c) 2026 Damien Flury, OST Ostschweizer Fachhochschule
%
\documentclass[tikz]{standalone}
\usepackage{amsmath}
\usepackage{times}
\usepackage{txfonts}
\usepackage{pgfplots}
\usetikzlibrary{arrows,math}
\definecolor{darkgreen}{rgb}{0,0.6,0}
\begin{document}
\def\skala{1}
\begin{tikzpicture}[>=latex,thick,scale=\skala,yscale=1.5]

% Flächen unter dem Integranden, Nullstelle bei x = -1
\fill[darkgreen!30,opacity=0.5]
	(-4,0) -- plot[domain=-4:-1,samples=100] ({\x},{(\x+1)/(1+\x*\x)}) -- cycle;
\fill[darkgreen!30,opacity=0.5]
	(-1,0) -- plot[domain=-1:4,samples=200] ({\x},{(\x+1)/(1+\x*\x)}) -- (4,0) -- cycle;

% Achsen
\draw[->] (-4.3,0) -- (4.5,0) coordinate[label={$x$}];
\draw[->] (0,-0.6) -- (0,1.5) coordinate[label={right:$y$}];

% Achsenbeschriftung
\foreach \x in {-4,-3,-2,-1,1,2,3,4}{
	\draw (\x,-0.033) -- (\x,0.033);
}
\foreach \x in {-4,-3,-2}{
	\node at (\x,0) [below] {$\x$};
}
\node at (-1,0) [below right] {$-1$};
\foreach \x in {1,2,3,4}{
	\node at (\x,0) [below] {$\x$};
}
\draw (-0.05,1) -- (0.05,1);
\node at (0,1) [above left] {$1$};

% Graph
\draw[color=darkgreen,line width=1.4pt]
	plot[domain=-4.2:4.2,samples=200] ({\x},{(\x+1)/(1+\x*\x)});

\node[color=darkgreen] at (3.0,0.85) {$f(x)=\dfrac{x+1}{1+x^2}$};

\end{tikzpicture}
\end{document}
